% ---------------------------------------------------------------------------
% Appendix for the tutor-as-diagnostician experiment (ACL style).
%
% Requires in the preamble:
%   \usepackage{booktabs}
%   \usepackage{listings}
%   \usepackage{xcolor}
%   \usepackage{amsmath}
%   \usepackage{multirow}
%
% and the prompt-listing style defined in Sec.~\ref{app:prompts} below
% (\lstdefinestyle{promptstyle}).
% ---------------------------------------------------------------------------

\appendix

\section{Experimental Setup}
\label{app:setup}

\subsection{Task Formulation}
\label{app:task}

A simulated student is conditioned on a ground-truth proficiency profile
$P_{\text{true}} : \mathcal{T} \rightarrow \{1,2,3,4,5\}$, where $\mathcal{T} =
\{S_1,\dots,S_{12}\}$ is the twelve-stage cognitive flow of
Table~\ref{tab:stages}. A tutor model interacts with the student for at most
$N$ turns and must recover $P_{\text{true}}$ without ever observing it. Each
turn the tutor emits a complete belief vector $P^{(t)}_{\text{pred}}$ over all
twelve stages; the belief at the terminal turn is taken as
$P_{\text{pred}}$.

The rollout proceeds as follows. At $t=0$ the student is shown the problem and
produces a first attempt (a fixed, non-adaptive prompt identical across all
conditions). For $t = 1,\dots,N$ the tutor observes the transcript so far and
emits a JSON action containing a natural-language probe, the stages that probe
targets, a belief vector, per-stage confidences, and a binary \texttt{stop}
flag. If \texttt{stop} is set the interview terminates early and that turn's
belief is final; otherwise the probe is sent to the student and the loop
continues.

Requiring a belief vector at \emph{every} turn (rather than only at the end)
is what makes convergence measurable: it yields an error-versus-turn curve
instead of a single endpoint. The explicit \texttt{stop} flag separates
``the tutor believes it is finished'' from ``the tutor exhausted its budget,''
which are different diagnoses of the same final error.

\subsection{Stage Taxonomy}
\label{app:stages}

\begin{table*}[t]
\centering
\small
\begin{tabular}{@{}llp{5.0cm}p{4.6cm}@{}}
\toprule
\textbf{Stage} & \textbf{Band} & \textbf{Captures} & \textbf{Representative bug class} \\
\midrule
$S_1$ Language Understanding  & pre-coding  & Parse statement, identify I/O and constraints & Inclusive bound read as exclusive \\
$S_2$ Problem Decomposition   & pre-coding  & Split into sub-problems, enumerate edge cases & Missing branch for empty input \\
$S_3$ Algorithm Formation     & pre-coding  & Devise a solution strategy & Greedy choice where DP is required \\
$S_4$ Algorithm Correctness   & pre-coding  & Verify the approach is sound & Recursion with an incorrect base case \\
$S_5$ Complexity Analysis     & pre-coding  & Reason about time/space cost & Nested scan in place of a hash lookup \\
$S_6$ Approach Finalization   & pre-coding  & Pick data structures, fix the plan & List used for membership tests \\
$S_7$ Code Translation        & code-level  & Turn the algorithm into code & Planned step omitted in implementation \\
$S_8$ Syntax Proficiency      & code-level  & Use language constructs correctly & Integer division for a real quantity \\
$S_9$ Code Correctness        & code-level  & Code matches intended logic & \texttt{<} in place of \texttt{<=} \\
$S_{10}$ Debugging            & code-level  & Spot and fix bugs on the fly & Shadowed variable, silent wrong value \\
$S_{11}$ Test-Case Checking   & code-level  & Validate with tests, edge cases & Passes visible tests, fails on duplicates \\
$S_{12}$ Verification         & code-level  & Confirm output, refine/optimize & Correct value, wrong return type \\
\bottomrule
\end{tabular}
\caption{The twelve-stage cognitive flow used as the latent skill space. The
tutor is shown the stage name, band, \emph{Captures} description and
representative bug class, together with the generic level scale of
Table~\ref{tab:levels}; it is never shown per-stage behavioral anchors.}
\label{tab:stages}
\end{table*}

\begin{table}[t]
\centering
\small
\begin{tabular}{@{}cp{5.6cm}@{}}
\toprule
\textbf{Level} & \textbf{Descriptor} \\
\midrule
1 & \textbf{Absent.} Does not perform this step at all, or performs it in a way that is actively wrong. \\
2 & \textbf{Fragile.} Attempts the step but usually gets it wrong; needs heavy scaffolding. \\
3 & \textbf{Inconsistent.} Correct on familiar or simple cases; breaks down as the problem varies. \\
4 & \textbf{Solid.} Usually correct; minor slips only under added complexity. \\
5 & \textbf{Fluent.} Reliably correct, self-checks unprompted, can explain \emph{why}. \\
\bottomrule
\end{tabular}
\caption{The proficiency scale, applied independently to each of the twelve
stages. The same descriptors are shown to both tutor and student.}
\label{tab:levels}
\end{table}

\subsection{Ground-Truth Profile Population}
\label{app:profiles}

Uniformly random profiles make a weak test set: they average out, and a
constant predictor scores well against them. We therefore sample profiles from
a mixture of archetypes (Table~\ref{tab:archetypes}), each targeting a distinct
failure mode of the tutor. Levels are jittered by $\pm 1$ with probability
$0.35$ (except \texttt{flat\_mid}) so that archetypes are not degenerate
constants. The population is deterministic given the seed, so every tutor is
evaluated against an identical set of students on identical problems.

For the $n{=}1000$ population the resulting level histogram is
$\{1{:}1438,\, 2{:}2430,\, 3{:}2513,\, 4{:}4005,\, 5{:}1614\}$ over
$1000 \times 12$ stage assignments, with $957$ distinct profile vectors. The
\texttt{spiky\_deficit} weak stage is cycled deterministically across
$S_1 \dots S_{12}$ rather than sampled, giving exactly $20$ single-stage
deficits per stage and hence even per-stage coverage. The $40$
\texttt{flat\_mid} profiles are deliberately identical, providing a
test--retest reliability estimate at no additional cost.

\begin{table}[t]
\centering
\small
\begin{tabular}{@{}lrp{3.5cm}@{}}
\toprule
\textbf{Archetype} & \textbf{$n$} & \textbf{Shape} \\
\midrule
\texttt{uniform\_random}        & 300 & unstructured \\
\texttt{spiky\_deficit}         & 240 & one stage at 1, rest $\approx 4$ \\
\texttt{band\_split\_plan\_weak}& 120 & $S_1$--$S_6$ low, $S_7$--$S_{12}$ high \\
\texttt{band\_split\_code\_weak}& 120 & $S_1$--$S_6$ high, $S_7$--$S_{12}$ low \\
\texttt{novice}                 & 100 & $\approx 2$ everywhere \\
\texttt{strong}                 &  80 & $\approx 4$ everywhere \\
\texttt{flat\_mid}              &  40 & exactly 3 everywhere \\
\midrule
\textbf{Total}                  & \textbf{1000} & \\
\bottomrule
\end{tabular}
\caption{Archetype mixture for the $n{=}1000$ profile population
(\texttt{profile\_seed}$=0$).}
\label{tab:archetypes}
\end{table}

\subsection{Conditions}
\label{app:conditions}

Localization error is not interpretable in isolation, so each tutor is
evaluated under three levels of access to the student:

\begin{itemize}
\itemsep0em
\item \textbf{\texttt{interactive}} --- the full protocol: the student's first
attempt followed by $N$ probe/response rounds ($\approx 2N{+}1$ model calls).
\item \textbf{\texttt{one\_shot}} --- the tutor observes only the student's
first attempt and must score immediately, with no probing (2 calls). The gap
to \texttt{interactive} isolates the \emph{value of probing}.
\item \textbf{\texttt{blind}} --- the tutor never observes the student at all
and predicts from the problem statement and priors over CS1 students alone
(1 call). This is a prior-only floor and doubles as a validity check: if
\texttt{interactive} does not improve on \texttt{blind}, either the simulated
student is not expressing its latent vector or the tutor is not using the
dialogue.
\end{itemize}

\subsection{Models}
\label{app:models}

The student model is held fixed at \texttt{openai/gpt-oss-20b} in every
condition and for every tutor, so that the task itself is constant across the
comparison. Table~\ref{tab:models} lists the tutors under evaluation. Note
that the \texttt{gpt-oss-20b} tutor cell is \emph{self-play}: tutor and
student share a distribution, so its score carries a shared-idiosyncrasy
advantage of unknown magnitude and is not directly comparable to the
cross-model cells. It is reported as a control rather than as a peer system.

\begin{table}[t]
\centering
\small
\begin{tabular}{@{}llr@{}}
\toprule
\textbf{Role} & \textbf{Model} & \textbf{Ctx.} \\
\midrule
Student (fixed) & \texttt{openai/gpt-oss-20b} & 16{,}384 \\
\midrule
\multirow{5}{*}{Tutor}
 & \texttt{meta-llama/Llama-3.1-8B-Instruct}    & 16{,}384 \\
 & \texttt{Qwen/Qwen2.5-Coder-14B-Instruct}     & 16{,}384 \\
 & \texttt{openai/gpt-oss-20b}$^{\dagger}$      & 16{,}384 \\
 & \texttt{google/gemma-4-31B-it}               & 12{,}288 \\
 & \texttt{Qwen/Qwen3-Coder-30B-A3B-Instruct}   & 12{,}288 \\
\bottomrule
\end{tabular}
\caption{Models. $^{\dagger}$self-play control (identical to the student
model). Context lengths are the vLLM \texttt{--max-model-len} used when
serving; the two largest models use a shorter context so that weights and KV
cache fit a single 80\,GB device.}
\label{tab:models}
\end{table}

\subsection{Decoding and Serving Configuration}
\label{app:config}

All models are served locally with vLLM 0.26.0 on NVIDIA A100-SXM4-80GB
devices (one device per role, tensor-parallel size 1,
\texttt{--gpu-memory-utilization} 0.90 for the student and 0.92 for the
tutor, server seed 42). Table~\ref{tab:hparams} lists the decoding and
protocol hyperparameters.

The tutor is decoded at a lower temperature than the student: the tutor must
emit parseable JSON and a stable belief, whereas the student is meant to
exhibit the variability of a real learner. For reasoning models exposing a
\texttt{reasoning\_effort} control, it is set to \texttt{low} independently per
role; the parameter is omitted entirely for models that do not implement it.
Because such models may return an empty \texttt{content} field with the
generation placed in a separate reasoning channel, the client falls back to
that channel when \texttt{content} is empty.

\begin{table}[t]
\centering
\small
\begin{tabular}{@{}lr@{}}
\toprule
\textbf{Parameter} & \textbf{Value} \\
\midrule
Probe budget $N$                    & 8 \\
Profiles                            & 1000 \\
Problems per profile                & 1 \\
Conditions                          & 3 \\
Episodes per tutor                  & 3000 \\
Model calls per tutor ($\approx$)   & 20{,}000 \\
\midrule
Tutor temperature                   & 0.3 \\
Student temperature                 & 0.7 \\
\texttt{max\_tokens} (both roles)   & 4096 \\
\texttt{reasoning\_effort}          & low$^{\ddagger}$ \\
Request timeout (s)                 & 180 \\
Client-side retries                 & 2 (exp.\ backoff) \\
Episode concurrency                 & 16 \\
\texttt{profile\_seed}              & 0 \\
\bottomrule
\end{tabular}
\caption{Protocol and decoding hyperparameters. $^{\ddagger}$applied only to
models implementing the parameter.}
\label{tab:hparams}
\end{table}

\subsection{Problem Source}
\label{app:problems}

Problems are drawn from the \texttt{test} split of
\texttt{newfacade/LeetCodeDataset}, filtered to \texttt{difficulty}
$=$ \texttt{Easy} ($48$ problems available). Each profile is interviewed on
the first $k$ problems of that filtered split, with $k$ fixed across tutors.
Easy competitive-programming problems are an approximation of CS1 coursework
rather than a substitute for it; we note this as a construct-validity
limitation in Section~\ref{app:limitations}.

\subsection{Metrics}
\label{app:metrics}

\paragraph{Localization accuracy.} The mean absolute distance between the
predicted and true profile,
\begin{equation}
d(P_{\text{pred}}, P_{\text{true}}) =
\frac{1}{|\mathcal{T}|} \sum_{c \in \mathcal{T}}
\bigl| P_{\text{pred}}(c) - P_{\text{true}}(c) \bigr| .
\end{equation}

Because a constant predictor already attains a low MAE against a realistic
profile distribution, we report it against the best constant baseline chosen
with hindsight over the same population,
$\text{MAE}^{*}_{\text{const}} = \min_{v \in \{1..5\}} \mathbb{E}\,
d(\mathbf{v}, P_{\text{true}})$, and define
\begin{equation}
\Delta_{\text{skill}} = \text{MAE}^{*}_{\text{const}} - d(P_{\text{pred}}, P_{\text{true}}),
\end{equation}
where $\Delta_{\text{skill}} \le 0$ indicates the tutor extracted no usable
information about the individual student.

\paragraph{Convergence cost.} Reported three ways, since they can disagree
informatively: (i) \emph{declared stop turn}, the turn at which the tutor set
\texttt{stop}; (ii) \emph{stabilization turn}, the first turn after which the
belief vector moves by at most $\epsilon = 0.25$ per stage for $2$ consecutive
turns; and (iii) \emph{turns to best}, the first turn at which the error curve
reaches its final value. We additionally report AUC-MAE, the mean of the
per-turn error curve, which rewards being both accurate and early.

\paragraph{Supporting diagnostics.} Signed bias (systematic leniency),
Spearman $\rho$ between predicted and true profile (shape recovery
independent of calibration), within-one-level rate, exact-match rate,
bottom-$3$ deficit recall, predicted level spread (range compression),
per-stage MAE, probe targeting counts, JSON parse-failure rate, and a
student-leakage rate that flags episodes in which the simulated student named
a stage, a rubric, or a numeric self-rating.

\subsection{Prompts}
\label{app:prompts}

\lstdefinestyle{promptstyle}{
  basicstyle=\ttfamily\scriptsize,
  breaklines=true,
  breakatwhitespace=false,
  columns=fullflexible,
  frame=single,
  framesep=4pt,
  xleftmargin=4pt,
  showstringspaces=false
}

\paragraph{Tutor system prompt.} \texttt{\{rubric\}} is substituted with
Tables~\ref{tab:stages} and~\ref{tab:levels} rendered as plain text, and
\texttt{\{budget\}} with $N$.

\begin{lstlisting}[style=promptstyle]
You are an expert CS1 instructor running a diagnostic interview with one
student about one problem. Your only goal is to work out *where in their
problem-solving process* they are weak -- not to teach, not to give them
the answer.

{rubric}

You get at most {budget} turns. On each turn you send the student one
message and see their reply. After each reply you must report your current
best estimate of their level (1-5) for all 12 stages, even the ones you
have no evidence about yet (guess, and refine later).

Probing advice: a probe that only asks "do you understand?" tells you
nothing. Good probes force the student to *perform* the stage you are
unsure about -- ask them to trace an input, predict an output, justify a
choice, produce a test case, find a bug you point at, or state a
complexity. You may only send messages a real tutor could send; do not ask
the student to reveal hidden state or rate themselves numerically.

Respond on EVERY turn with a single JSON object and nothing else:

{
  "reasoning": "<2-4 sentences: what the last reply told you and what you still need>",
  "target_stages": ["S3", "S5"],
  "probe": "<the exact message to send the student>",
  "belief": {"S1": 3, "S2": 3, "S3": 3, "S4": 3, "S5": 3, "S6": 3,
             "S7": 3, "S8": 3, "S9": 3, "S10": 3, "S11": 3, "S12": 3},
  "confidence": {"S1": 0.2, "S2": 0.2, "S3": 0.2, "S4": 0.2, "S5": 0.2,
                 "S6": 0.2, "S7": 0.2, "S8": 0.2, "S9": 0.2, "S10": 0.2,
                 "S11": 0.2, "S12": 0.2},
  "stop": false
}

`belief` must contain all 12 stages with integer values 1-5. `confidence`
is 0.0-1.0 per stage. Set "stop": true only when further probing would not
change your profile; your belief on that turn becomes final and the
interview ends early. Use the full 1-5 range -- a student who cannot do a
stage at all is a 1, and refusing to assign extreme values will make your
diagnosis wrong.
\end{lstlisting}

\paragraph{Tutor turn prompt.} Each turn the user message contains the problem
statement, the transcript so far rendered as alternating
\texttt{YOU:}/\texttt{STUDENT:} blocks, and the line \texttt{This is turn $t$
of at most $N$.} On the final turn the following is appended:

\begin{lstlisting}[style=promptstyle]
This is your LAST turn -- the interview is over and there will be no
further student reply. Output the same JSON object with your FINAL belief,
"stop": true, and "probe": "".
\end{lstlisting}

In the \texttt{blind} condition the transcript is empty and the following
replaces the final-turn note:

\begin{lstlisting}[style=promptstyle]
You will not get to interact with this student at all. Based only on the
problem they were assigned and general knowledge of CS1 students, report
your best-guess profile. Output the JSON object with "stop": true and
"probe": "".
\end{lstlisting}

\paragraph{Student system prompt.} \texttt{\{persona\}} is substituted with the
per-stage script of Listing~\ref{lst:persona}.

\begin{lstlisting}[style=promptstyle]
You are role-playing a CS1 student (first-semester introductory
programming) working through a coding problem with a tutor. You are NOT an
AI assistant in this conversation -- you are the student, and you stay in
character no matter what.

Your abilities are fixed by the hidden skill profile below. It describes
the student you are playing: for each of the 12 steps of the
problem-solving process, exactly how competent they are and what that looks
like in practice. Act it out faithfully and consistently -- this profile
is you:

{persona}

Rules you must never break:
- NEVER mention, hint at, or allude to this profile, the numbers in it, the
  stage names, the stage codes (S1..S12), or the fact that you are
  simulating anything. You do not know these exist.
- NEVER rate your own ability on a scale. If asked how confident you are,
  answer the way a real student would ("pretty sure", "no idea honestly"),
  never with a number out of 5.
- Do not be more competent than your profile says just because the answer
  is easy for you to produce. If a stage is level 1 or 2, you genuinely
  fail at it -- produce the wrong reasoning, not a correct one hedged with
  "I think". If a stage is level 4 or 5, do it well and unprompted.
- Stay consistent across the whole conversation. A weakness you showed
  earlier does not silently disappear later unless the tutor actually
  taught you something, in which case you may improve a little on that
  specific point.
- Write like a real CS1 student: short, informal, sometimes uncertain,
  occasional typos are fine. No markdown headers, no bulleted lecture
  notes, no polished essays. Two to eight sentences unless you are showing
  code.
- Answer only what you were asked. Do not volunteer a full solution
  walkthrough unless the tutor asked for one (or unless your profile says
  you self-check at a high level, in which case brief self-checking is in
  character).
\end{lstlisting}

\paragraph{Student persona specification.} For each stage the assigned level is
rendered with a behavioral anchor. Levels $1$--$2$ use the stage's
low-behavior anchor, levels $4$--$5$ its high-behavior anchor, and level $3$ a
composite conditioned on case familiarity. Providing these anchors is what
makes a latent level observable in dialogue; without them the student model
regresses toward a generic ``average student'' and the latent signal is
destroyed.

\begin{lstlisting}[style=promptstyle,caption={Persona entry format, shown for one stage.},label={lst:persona}]
S5 - Complexity Analysis => LEVEL 2 (Fragile).
    This step covers: Reason about time and space cost.
    This student: Guesses big-O, or answers in vague terms ('fast', 'not
    too slow'); miscounts nested loops; ignores space entirely.
\end{lstlisting}

\paragraph{Student first-turn prompt.} Identical in every episode and
condition, so that the first attempt is not adapted to the tutor.

\begin{lstlisting}[style=promptstyle]
Here's the problem you've been assigned:

{problem}

Take a first pass at it. Briefly say how you're reading the problem and
what your plan is, then write your attempt at the code in a ```python
block. Do this the way you naturally would -- don't try to be thorough if
that's not how you work.
\end{lstlisting}

\subsection{Robustness and Implementation Details}
\label{app:robustness}

\paragraph{Output parsing.} Tutor actions are extracted with a tolerant parser
that accepts a fenced JSON block, the widest brace-delimited span, or a bare
object, and coerces out-of-range levels into $\{1,\dots,5\}$. A malformed
response is retried once with the parser's complaint appended; if it still
fails, the previous turn's belief is carried forward and the turn is recorded
as a format failure. Parse-failure rate is reported as a compliance metric
rather than silently discarded.

\paragraph{Chat-template compatibility.} Some chat templates do not define a
\texttt{system} role. When a request is rejected for this reason, the client
folds the system prompt into the first user message and retries, so that
template differences do not confound the model comparison.

\paragraph{Episode independence.} Each episode is an independent conversation;
no state carries across episodes. The student is re-prompted with the full
transcript each turn so that its behavior remains consistent within an
episode.

\subsection{Limitations}
\label{app:limitations}

Every measurement here is conditional on the simulated student faithfully
expressing $P_{\text{true}}$. The \texttt{blind} condition and the
student-leakage rate are built-in validity checks, but neither establishes
fidelity; they only detect gross failure. A stronger check, which we do not
perform, is a paired design over profiles differing in exactly one stage,
verifying that the tutor's belief moves on that stage and not on others.

A capable student model may also be unable to convincingly portray level-1
behavior, since competence leaks through; error at the low end of the scale
should therefore be interpreted with care and compared across student models
before being attributed to the tutor. Finally, the twelve stages are scored
independently although they are not independent in practice (a weak $S_3$
constrains $S_7$), and mean absolute error ignores this structure.
